\chapter{Proofs of the SF axioms for the braiding functor data}

\section{Overview}

In Chapter 6 the braiding functor data was described by generators
$\beta_0, \beta_1, \beta_2$ subject to the truncated axioms
$[\mathrm{SF}_1]$--$[\mathrm{SF}_4]$. In the formalization this description is
now a \emph{theorem} rather than a definition: the data is defined to be an
\Ainf-functor
\[
  F \colon \KLRW \longrightarrow K^\bullet\big(\Add(\KLRW)\big)
\]
whose components vanish above level two, and the four axioms are derived by
evaluating the general \Ainf-functor equation $[\mathrm{SF}_n]$ on chains of
degree-zero morphisms. The logical shape of each derivation is important: the
equation $[\mathrm{SF}_n]$ --- ``insertion side $=$ partition side'' --- holds
\emph{by definition} of the datum, and the proof consists of computing what
its two sides are. The resulting identities are therefore not structural
facts about the target category; they are the constraints that the functor
equation imposes on the components $\beta_1, \beta_2$. This chapter records
the mathematics of these computations, in the exact shape mirrored by the
Lean proofs.

\paragraph{Convention.} Composition is written diagrammatically by
juxtaposition: for $f \in \Hom(A,B)$ and $g \in \Hom(B,C)$ we write
$fg \in \Hom(A,C)$ (Lean's $f \gg g$). Chapter 6 writes the same composite as
$g \circ f$; the axioms below agree with the ones stated there after this
translation.

\section{The two \texorpdfstring{\Ainf}{A-infinity}-categories}

\begin{definition}[$\KLRW$ as a graded category]\label{def:klrw-graded}
  For $A, B \in \KLRW$ put
  \[
    \Hom_{\KLRW}(A,B)_d \;=\;
    \begin{cases}
      \StrandSpace(R) & d = 0,\\
      0 & d \neq 0.
    \end{cases}
  \]
  The \Ainf-operations are: $\mu^{\KLRW}_2$ on a chain of degrees $(0,0)$ is
  the convolution composition of Chapter 6, and every other operation is zero
  ($\mu^{\KLRW}_1 = 0$, $\mu^{\KLRW}_n = 0$ for $n \geq 3$, and $\mu^{\KLRW}_2$
  vanishes on chains whose degree vector is not $(0,0)$ --- where some argument
  module is the zero module anyway).
\end{definition}

\begin{remark}[the grading forces the degeneracy]
  The operation $\mu_n$ has target degree $\sum_i d_i + (2-n)$. On a chain of
  degree-zero morphisms this equals $2-n$, which is nonzero unless $n = 2$; so
  any multilinear $\mu_n$ with $n \neq 2$ takes values in the zero module and
  the degenerate structure of Definition~\ref{def:klrw-graded} is the only
  possible one. (Lean: \texttt{klrwHom}, \texttt{klrwHom\_eq\_zero},
  \texttt{klrwAInfinityPreCategory}.)
\end{remark}

\begin{theorem}[$\KLRW$ is an \Ainf-category]\label{thm:klrw-ainf}
  The operations of Definition~\ref{def:klrw-graded} satisfy the Stasheff
  identities.
\end{theorem}

\begin{proof}
  The $n$-th Stasheff identity is a signed sum of terms
  $\mu_{n+1-s}(\dots, \mu_s(\dots), \dots)$. If $s \neq 2$ the inner operation
  vanishes; if $n + 1 - s \neq 2$ the outer one does. Both equal $2$ only when
  $n = 3$ and $s = 2$, so every identity except the third is a sum of zeros.
  For $n = 3$, if some argument has nonzero degree it lies in the zero module
  and all terms vanish; otherwise the two surviving insertion terms are
  $\mu_2(\mu_2(x_0,x_1),x_2)$ and $\mu_2(x_0,\mu_2(x_1,x_2))$, which carry
  opposite signs and cancel by associativity of convolution
  (Chapter 6, via $\operatorname{compShift}$-associativity). (Lean:
  \texttt{klrwAInfinityCategory} in \texttt{KLRW.lean}.)
\end{proof}

\begin{definition}[the dg target]\label{def:dg-target}
  Throughout this chapter we abbreviate the target category by
  \[
    B \;:=\; K^\bullet\big(\Add(\KLRW)\big),
  \]
  and write $\mu^B_n$ for its operations (matching the notation of
  Chapter 6). It carries the \Ainf-structure of a dg-category:
  \[
    \mu^B_1 = \delta \quad\text{(the Hom-complex differential)}, \qquad
    \mu^B_2(x, y) = (-1)^{|y|}\, x y, \qquad
    \mu^B_n = 0 \ (n \geq 3),
  \]
  where $|y|$ is the degree of $y$ and $xy$ is composition of Hom-complex
  elements. The Stasheff identities are $\delta^2 = 0$, the Leibniz rule, and
  strict associativity. (Lean: \texttt{bccAInfinityPreCategory} and the
  \texttt{AInfinityCategory} instance in \texttt{BoundedCochainComplex.lean}.)
\end{definition}

\section{\texorpdfstring{\Ainf}{A-infinity}-functors and the truncation}

\begin{definition}[\Ainf-functor]\label{def:ainf-functor}
  An \Ainf-functor $F \colon \mathcal{A} \to \mathcal{B}$ consists of an
  object map $F_0$ and, for each $n \geq 1$, multilinear components $f_n$ on
  composable chains $(a_1, \dots, a_n)$, of target degree
  $\sum_i |a_i| + (1 - n)$, satisfying for every $n$ and every chain the
  equation $[\mathrm{SF}_n]$, in degree $\sum_i |a_i| + (2-n)$:
  \begin{multline*}
    \sum_{\substack{r + s + t = n\\ s \geq 1}}
      \varepsilon_{r,s}\,
      f_{r+1+t}\big(a_1, \dots, a_r,\ \mu^{\mathcal A}_s(a_{r+1}, \dots, a_{r+s}),\
        a_{r+s+1}, \dots, a_n\big) \\
    = \sum_{\substack{(c_1, \dots, c_k) \vDash n}}
      \mu^{\mathcal B}_k\big(f_{c_1}(a_1, \dots, a_{c_1}),\ \dots,\
        f_{c_k}(a_{n-c_k+1}, \dots, a_n)\big),
  \end{multline*}
  where the right sum runs over all compositions of $n$ (ordered tuples of
  positive integers summing to $n$). Note the asymmetry: the insertion side
  applies operations of the \emph{source} $\mathcal{A}$ inside the arguments
  of a single component, while the partition side applies one operation of
  the \emph{target} $\mathcal{B}$ to a tuple of component values. The
  insertion sign is
  \[
    \varepsilon_{r,s} \;=\; (-1)^{\,|a_{r+s+1}| + \dots + |a_n| \, +\, t},
    \qquad t = n - r - s .
  \]
  (Lean: \texttt{AInfinityFunctor}, with the two sides
  \texttt{indexedSFLeftSum} and \texttt{indexedSFRightSum} in
  \texttt{AInfinityFunctor.lean}.)
\end{definition}

\begin{remark}[sign convention]
  The insertion sign follows the same convention as the formalized Stasheff
  sum; the partition side carries no additional sign. Other references
  distribute signs differently; over a base ring with $2 = 0$ --- the standing
  assumption for the braiding data --- all conventions agree.
\end{remark}

\begin{definition}[braiding functor data]\label{def:bfd}
  A \emph{braiding functor datum} is an \Ainf-functor
  $F \colon \KLRW \to K^\bullet(\Add(\KLRW))$ together with the truncation
  condition $f_k = 0$ for all $k \geq 3$. From it one recovers the generators
  of Chapter 6:
  \[
    \beta_0 := F_0, \qquad
    \beta_1(f) := f_1(f) \in \Hom^{0}\big(\beta_0 A, \beta_0 B\big), \qquad
    \beta_2(f,g) := f_2(f,g) \in \Hom^{-1}\big(\beta_0 A, \beta_0 C\big),
  \]
  the degrees being $1 - 1 = 0$ and $1 - 2 = -1$ since all inputs sit in
  degree $0$. (Lean: \texttt{BraidingFunctorData} with fields
  \texttt{toFunctor}, \texttt{trunc}, and derived \texttt{gen₀},
  \texttt{gen₁}, \texttt{gen₂}.)
\end{definition}

Throughout the rest of the chapter, $f, g, h, k$ denote composable morphisms
of $\KLRW$; all have degree $0$, so every insertion sign reduces to
\[
  \varepsilon_{r,s} = (-1)^{t} = (-1)^{\,n-r-s}.
\]

\section{Vanishing criteria}

\begin{lemma}[term vanishing]\label{lem:vanishing}
  Fix a chain of length $n$.
  \begin{enumerate}
    \item If the outer component $f_{n+1-s}$ vanishes identically, the
      $(r,s)$ insertion term is zero.
    \item If the inner operation $\mu^{\mathcal A}_s$ vanishes identically,
      the $(r,s)$ insertion term is zero.
    \item If $\mu^{\mathcal B}_k$ vanishes identically, the partition term at
      any composition of length $k$ is zero.
    \item If $f_{c_j}$ vanishes identically for some block $c_j$ of the
      composition $c$, the partition term at $c$ is zero.
  \end{enumerate}
\end{lemma}

\begin{proof}
  (1) and (3) are immediate. For (2) and (4), the vanishing operation makes
  one coordinate of the input vector of the outer map zero, and a multilinear
  map vanishes on any vector with a zero coordinate. (Lean:
  \texttt{indexedSFLeftTerm\_outer\_zero},
  \texttt{indexedSFLeftTerm\_inner\_zero},
  \texttt{indexedSFRightTerm\_outer\_zero},
  \texttt{indexedSFRightTerm\_inner\_zero}.)
\end{proof}

\section{Normal forms and evaluation of the sums}\label{sec:normal-forms}

On paper one discards the vanishing terms and reads off the identity. In the
formalization two further kinds of manipulation are needed, because the terms
produced by the general $[\mathrm{SF}_n]$ machinery are not
\emph{syntactically} the expressions appearing in the axioms. We record them
once here; the proofs of Theorems~\ref{thm:sf1}--\ref{thm:sf4} then read as
ordinary calculations, and use these facts silently at every step.

\begin{remark}[canonical identifications]\label{rem:canonical}
  The framework assigns to each operation a target degree that is a
  \emph{term}, e.g.\ $\sum_i d_i + (2-n)$, not the integer it denotes, and
  moves values between the corresponding modules by transports. On a chain of
  literal degree-zero morphisms every degree expression in sight is closed
  and evaluates, so all these transports are along equations between
  definitionally equal integers and collapse (proof irrelevance). Formally
  one strips them by passing to heterogeneous equality and converts back at
  the end, both sides then lying in the same degree. In the proofs below we
  never mention them again.
\end{remark}

\begin{lemma}[normal form of a term]\label{lem:normal-form}
  On a degree-zero chain, each insertion term equals the expression
  \[
    f_{n+1-s}\big(a_1, \dots, a_r,\
      \mu^{\KLRW}_s(a_{r+1}, \dots, a_{r+s}),\ a_{r+s+1}, \dots, a_n\big)
  \]
  it denotes, evaluated on the collapsed chain in its literal degrees; and
  each partition term equals $\mu^B_k$ applied to the block evaluations
  $f_{c_j}(\dots)$. Moreover $\mu^{\KLRW}_2(f,g) = fg$ holds
  \emph{definitionally}, so composites may be substituted freely.
\end{lemma}

\begin{proof}
  This is where the formal work sits, and it is not purely definitional. The
  machine-built term is a transport of the component applied to a
  \emph{re-indexed} chain (the inner window collapsed, or a block cut out).
  The re-indexed object, degree and morphism assignments agree with the
  displayed ones at every individual index --- each instance a definitional
  identity, since the index arithmetic evaluates there --- but not as
  functions: at a variable index the bodies differ, e.g.\ $0 + i$ and $i$
  are distinct normal forms in $\N$. So the two applications of the
  (multilinear) component are equated by congruence: function extensionality
  for the object and degree assignments, checked index by index; a
  heterogeneous function extensionality for the chains of morphisms, whose
  entries become definitional identities once their transports are stripped
  (Remark~\ref{rem:canonical}). For the last claim, the case distinction in
  the definition of $\mu^{\KLRW}_2$ is decided by evaluation on literal
  degrees, its internal casts are along $0 = 0$, and the convolution
  composite remains. (Lean: the congruence bridges built from
  \texttt{mlm\_apply\_congr\_heq} and \texttt{Function.hfunext} inside
  \texttt{sf}$_1$--\texttt{sf}$_4$.)
\end{proof}

\begin{lemma}[sum evaluation]\label{lem:sum-eval}
  On a degree-zero chain, each side of $[\mathrm{SF}_n]$ equals the sum of
  its non-vanishing terms, the insertion terms weighted by
  $(-1)^{\,n-r-s}$. Formally: vanishing summands are rewritten to zero by
  Lemma~\ref{lem:vanishing} and absorbed; the index sets are enumerated ---
  compositions of a small $n$ are classified by destructing the block list;
  and the signs, being closed expressions, are evaluated. (Lean:
  \texttt{Finset.sum\_eq\_single\_of\_mem} and a two-survivor variant;
  \texttt{composition\_two\_cases}, \texttt{composition\_three\_cases},
  \texttt{composition\_four\_eq\_two\_two}.)
\end{lemma}

\section{Arity 1: \texorpdfstring{$[\mathrm{SF}_1]$}{[SF1]}}

\begin{theorem}\label{thm:sf1}
  For every $f \in \Hom_{\KLRW}(A, B)$,
  \[
    \mu^B_1\big(\beta_1(f)\big) = 0 ,
  \]
  an identity in degree $1$; that is, $\beta_1(f)$ is a cycle for the
  Hom-complex differential --- a chain map.
\end{theorem}

\begin{proof}
  Since $F$ is an \Ainf-functor (Definition~\ref{def:bfd}), the equation
  $[\mathrm{SF}_1]$ of Definition~\ref{def:ainf-functor} holds on the chain
  $(f)$:
  \[
    \text{(insertion side at $(f)$)} \;=\; \text{(partition side at $(f)$)} .
  \]
  We compute the two sides. The insertion side applies source operations: its
  only term is $f_1\big(\mu^{\KLRW}_1 f\big)$, which vanishes since
  $\mu^{\KLRW}_1 = 0$ (Lemma~\ref{lem:vanishing}); so the insertion side is
  $0$. The partition side applies a target operation to component values: the
  only composition of $1$ is $(1)$, with term
  $\mu^{B}_1\big(f_1(f)\big) = \mu^{B}_1\big(\beta_1(f)\big)$ --- the
  Hom-complex differential of $B = K^\bullet(\Add(\KLRW))$ applied to
  $\beta_1(f)$; so the partition side is $\mu^{B}_1\big(\beta_1(f)\big)$.
  Substituting both computations into the displayed equation gives
  $0 = \mu^{B}_1\big(\beta_1(f)\big)$.
  (Lean: \texttt{BraidingFunctorData.sf}$_1$.)
\end{proof}

\section{Arity 2: \texorpdfstring{$[\mathrm{SF}_2]$}{[SF2]}}

\begin{theorem}\label{thm:sf2}
  For composable $f, g$,
  \[
    \beta_1(fg) \;=\; \mu^B_2\big(\beta_1(f), \beta_1(g)\big)
      \;+\; \mu^B_1\big(\beta_2(f,g)\big),
  \]
  an identity in degree $0$.
\end{theorem}

\begin{proof}
  As $F$ is an \Ainf-functor, $[\mathrm{SF}_2]$ holds on the chain $(f,g)$:
  the insertion side and the partition side at $(f,g)$ are equal. We compute
  the two sides.

  \emph{Insertion side.} Inserting $\mu^{\KLRW}_1$ into either argument gives
  $f_2(\mu^{\KLRW}_1 f,\, g)$ or $f_2(f,\, \mu^{\KLRW}_1 g)$, both zero since
  $\mu^{\KLRW}_1 = 0$. The remaining insertion is that of $\mu^{\KLRW}_2$,
  with sign $(-1)^{\,2-0-2} = +1$; since $\mu^{\KLRW}_2(f,g) = fg$, the
  insertion side equals
  \[
    f_1\big(\mu^{\KLRW}_2(f,g)\big) \;=\; f_1(fg) \;=\; \beta_1(fg).
  \]

  \emph{Partition side.} The compositions of $2$ are $(1,1)$ and $(2)$, and
  neither term vanishes; the side equals
  \[
    \mu^B_2\big(f_1(f),\, f_1(g)\big) + \mu^B_1\big(f_2(f,g)\big)
    \;=\; \mu^B_2\big(\beta_1(f), \beta_1(g)\big)
      + \mu^B_1\big(\beta_2(f,g)\big).
  \]

  Substituting the two computations into the defining equation gives the
  statement. (Lean: \texttt{BraidingFunctorData.sf}$_2$.)
\end{proof}

\section{Arity 3: \texorpdfstring{$[\mathrm{SF}_3]$}{[SF3]}}

\begin{theorem}\label{thm:sf3}
  For composable $f, g, h$,
  \[
    -\,\beta_2(fg,\, h) \;+\; \beta_2(f,\, gh)
    \;=\; \mu^B_2\big(\beta_1(f), \beta_2(g,h)\big)
      \;+\; \mu^B_2\big(\beta_2(f,g), \beta_1(h)\big),
  \]
  an identity in degree $-1$. Over a base with $2 = 0$ the sign disappears:
  \[
    \beta_2(f,\, gh) + \beta_2(fg,\, h)
    = \mu^B_2\big(\beta_1(f), \beta_2(g,h)\big)
      + \mu^B_2\big(\beta_2(f,g), \beta_1(h)\big).
  \]
\end{theorem}

\begin{proof}
  As $F$ is an \Ainf-functor, $[\mathrm{SF}_3]$ holds on the chain
  $(f,g,h)$: the two sides below are equal. We compute them.

  \emph{Insertion side.} Inserting $\mu^{\KLRW}_1$ (three ways) or
  $\mu^{\KLRW}_3$ (one way) gives zero. The two insertions of
  $\mu^{\KLRW}_2$ survive, with signs $(-1)^{\,3-r-2}$ for $r = 0, 1$:
  \[
    (-1)\cdot f_2\big(\mu^{\KLRW}_2(f,g),\, h\big)
      \;+\; f_2\big(f,\, \mu^{\KLRW}_2(g,h)\big)
    \;=\; -\,\beta_2(fg,\, h) + \beta_2(f,\, gh).
  \]

  \emph{Partition side.} Of the four compositions of $3$, the term at
  $(1,1,1)$ vanishes since $\mu^B_3 = 0$, and the term at $(3)$ vanishes
  since $f_3 = 0$ by truncation. The two survivors give
  \[
    \mu^B_2\big(f_1(f),\, f_2(g,h)\big)
      + \mu^B_2\big(f_2(f,g),\, f_1(h)\big)
    \;=\; \mu^B_2\big(\beta_1(f), \beta_2(g,h)\big)
      + \mu^B_2\big(\beta_2(f,g), \beta_1(h)\big).
  \]

  Substituting the two computations into the defining equation gives the
  signed identity. When $2 = 0$ the sign
  disappears, since $-\beta_2(fg, h) = \beta_2(fg, h)$ by the $2$-torsion of
  the Hom-groups (Lemma~\ref{lem:two-torsion}); this is the only point in the
  four proofs where the characteristic is used. (Lean:
  \texttt{BraidingFunctorData.sf}$_3$; the char-$2$ step is
  \texttt{sf₃\_key2}.)
\end{proof}

\section{Arity 4: \texorpdfstring{$[\mathrm{SF}_4]$}{[SF4]}}

\begin{theorem}\label{thm:sf4}
  For composable $f, g, h, k$,
  \[
    0 \;=\; \mu^B_2\big(\beta_2(f,g),\, \beta_2(h,k)\big),
  \]
  an identity in degree $-2$.
\end{theorem}

\begin{proof}
  As $F$ is an \Ainf-functor, $[\mathrm{SF}_4]$ holds on the chain
  $(f,g,h,k)$:
  \[
    \text{(insertion side at $(f,g,h,k)$)}
      \;=\; \text{(partition side at $(f,g,h,k)$)} .
  \]
  We compute the two sides.

  \emph{Insertion side.} Every term vanishes: inserting $\mu^{\KLRW}_1$,
  $\mu^{\KLRW}_3$ or $\mu^{\KLRW}_4$ gives zero because those operations
  vanish, and inserting $\mu^{\KLRW}_2$ leaves an outer component
  $f_{4+1-2} = f_3 = 0$ by truncation.

  \emph{Partition side.} Let $c$ be a composition of $4$. If $c$ has three or
  more blocks, its term vanishes since $\mu^B_k = 0$ for $k \geq 3$; if some
  block has size $\geq 3$, it vanishes by truncation. The only composition
  with at most two blocks, each of size at most $2$, summing to $4$, is
  $(2,2)$, whose term is
  \[
    \mu^B_2\big(f_2(f,g),\, f_2(h,k)\big)
    \;=\; \mu^B_2\big(\beta_2(f,g),\, \beta_2(h,k)\big).
  \]

  So the insertion side is $0$ and the partition side is
  $\mu^B_2\big(\beta_2(f,g), \beta_2(h,k)\big)$; substituting both into the
  defining equation gives
  \[
    0 \;=\; \mu^B_2\big(\beta_2(f,g),\, \beta_2(h,k)\big) .
  \]
  Note that this vanishing is not an identity of the target category ---
  $\mu^B_2$ certainly does not kill arbitrary pairs of degree-$(-1)$
  elements. It is a constraint that the functor equation places on the
  component $\beta_2$ of the datum. (Lean:
  \texttt{BraidingFunctorData.sf}$_4$; the classification of compositions is
  \texttt{composition\_four\_eq\_two\_two}.)
\end{proof}

\section{Arities \texorpdfstring{$\geq 5$}{>= 5}}

\begin{proposition}\label{prop:sf-high}
  For $n \geq 5$, both sides of $[\mathrm{SF}_n]$ vanish identically on
  degree-zero chains, for any truncated datum.
\end{proposition}

\begin{proof}
  Insertion side: if $s \leq 2$ the outer component has arity
  $n + 1 - s \geq n - 1 \geq 4$, hence vanishes by truncation; if $s \geq 3$
  the inner $\mu^{\KLRW}_s$ vanishes. Partition side: if some block is
  $\geq 3$ its component vanishes by truncation; if all blocks are $\leq 2$
  the length is at least $\lceil n/2 \rceil \geq 3$ and $\mu^B_k = 0$.
\end{proof}

\begin{remark}
  Proposition~\ref{prop:sf-high} is the reason the axiom list of Chapter 6 is
  finite. In the formalization it is the obligation discharged (at high
  arities) when \emph{constructing} a \texttt{BraidingFunctorData}; the
  four theorems above are the converse direction, extracting the axioms from
  the structure.
\end{remark}

\section{Translation to differential and composition}

\begin{lemma}[2-torsion]\label{lem:two-torsion}
  Assume $2 = 0$ in $R$. Then every element $w$ of every Hom-group of
  $K^\bullet(\Add(\KLRW))$ satisfies $-w = w$.
\end{lemma}

\begin{proof}
  Componentwise, such an element is a matrix of elements of
  $\StrandSpace(R) = \bigoplus_{\N} R$, and $-r = r$ in $R$ since $2 = 0$.
  (Lean: \texttt{finCochain\_neg\_eq\_self}.)
\end{proof}

Unfolding $\mu^B_1 = \delta$ and $\mu^B_2(x,y) = (-1)^{|y|} xy$ in
Theorems~\ref{thm:sf1}--\ref{thm:sf4} gives the axioms in the language of
Chapter 6 (writing $\cdot$ for composition of Hom-complex elements):
\begin{itemize}
  \item $[\mathrm{SF}_1]$: $\delta\big(\beta_1(f)\big) = 0$. No sign appears.
  \item $[\mathrm{SF}_2]$: $\beta_1(fg) = \beta_1(f)\cdot\beta_1(g)
    + \delta\big(\beta_2(f,g)\big)$, since the Koszul sign of $\mu^B_2$ at
    degrees $(0,0)$ is $(-1)^0 = +1$. Valid over any base ring.
  \item $[\mathrm{SF}_3]$: here $\mu^B_2\big(\beta_1(f), \beta_2(g,h)\big)
    = (-1)^{-1}\, \beta_1(f)\cdot\beta_2(g,h)$ carries a genuine sign
    ($|\beta_2| = -1$), while the second term does not
    ($|\beta_1| = 0$). Over $2 = 0$ all signs are absorbed by
    Lemma~\ref{lem:two-torsion}:
    \[
      \beta_2(f, gh) + \beta_2(fg, h)
      = \beta_1(f)\cdot\beta_2(g,h) + \beta_2(f,g)\cdot\beta_1(h).
    \]
    This is the only one of the four axioms whose composition-language form
    genuinely uses the characteristic.
  \item $[\mathrm{SF}_4]$: $\beta_2(f,g)\cdot\beta_2(h,k) = 0$; the Koszul
    sign $(-1)^{-1}$ is absorbed by the vanishing right-hand side, so this
    too is valid over any base ring.
\end{itemize}
(Lean: \texttt{sf}$_1$\texttt{\_fin}--\texttt{sf}$_4$\texttt{\_fin} in
\texttt{Braiding.lean}.)

\section{Correspondence with the formalization}

\begin{center}
\begin{tabular}{ll}
  \textbf{Mathematics} & \textbf{Lean} \\[2pt]
  $\KLRW$ as \Ainf-category (Thm.~\ref{thm:klrw-ainf}) &
    \texttt{klrwAInfinityPreCategory}, \texttt{klrwAInfinityCategory} \\
  dg structure on $K^\bullet(\Add(\KLRW))$ (Def.~\ref{def:dg-target}) &
    \texttt{bccAInfinityPreCategory} + instance \\
  \Ainf-functor, $[\mathrm{SF}_n]$ (Def.~\ref{def:ainf-functor}) &
    \texttt{AInfinityFunctor}, \texttt{indexedSFLeftSum},
    \texttt{indexedSFRightSum} \\
  truncated datum (Def.~\ref{def:bfd}) & \texttt{BraidingFunctorData} \\
  vanishing criteria (Lem.~\ref{lem:vanishing}) &
    \texttt{indexedSF*Term\_outer/inner\_zero} \\
  Theorems~\ref{thm:sf1}--\ref{thm:sf4} &
    \texttt{BraidingFunctorData.sf}$_1$--\texttt{sf}$_4$ \\
  composition-language forms &
    \texttt{sf}$_1$\texttt{\_fin}--\texttt{sf}$_4$\texttt{\_fin} \\
  2-torsion (Lem.~\ref{lem:two-torsion}) & \texttt{finCochain\_neg\_eq\_self} \\
\end{tabular}
\end{center}
